
\subsubsection{3.1 Quality Metrics for Attribution}

Two dimensions of attribution quality are formalized, each corresponding to different downstream applications.

\textbf{Rank Preservation (rho\_r).} Rank preservation measures how well the method preserves the ordering of influential examples relative to ground truth:

rho\_r = Spearman(phi\_hat, phi\_star)

where phi\_hat is the estimated influence vector and phi\_star is the ground truth. High rho\_r matters when practitioners need to identify the top-k most influential examples.

\textbf{Magnitude Fidelity (rho\_m).} Magnitude fidelity measures how accurately the method estimates influence magnitudes:

rho\_m = Pearson(phi\_hat, phi\_star)

where Pearson correlation captures both ranking and scale fidelity. High rho\_m matters for data valuation applications where practitioners need accurate influence values.

\subsubsection{3.2 Compute Normalization via Gradient-Equivalent Operations}

Direct comparison across methods is complicated by different computational primitives---TRAK uses random projections, IF uses HVP iterations, TracIn uses gradient similarities. Computation is normalized via gradient-equivalent operations (GEOs):

\textbf{Definition:} One GEO equals the cost of computing one gradient of the training loss with respect to model parameters.

\begin{table}[htbp]
\centering
\resizebox{\columnwidth}{!}{%
\begin{tabular}{ll}
\toprule
Method & Compute Budget Translation \\
\midrule
TRAK & GEOs = projection\_dim x samples\_per\_dim \\
IF & GEOs = hvp\_iterations x convergence\_checks \\
TracIn & GEOs = checkpoints x gradient\_computations \\
FastIF & GEOs = arnoldi\_iterations \\
\bottomrule
\end{tabular}
}
\end{table}

\subsubsection{3.3 Hypothesis-Driven Validation Structure}

The investigation is structured as four sub-hypotheses:

\textbf{H-E1: Existence of Pareto Trade-offs (MUST\_WORK).} At least one method pair exhibits statistically significant metric crossings with non-overlapping 95\% bootstrap confidence intervals at two or more compute levels. Gate: at least 2 budget levels with CI-separated crossings.

\textbf{H-M1: Convex Baseline Coupling (MUST\_WORK).} In convex settings, cross-metric partial correlations corr(rho\_r, rho\_m | b) at least 0.95 at all compute levels. Gate: minimum correlation at least 0.95 across all budgets.

\textbf{H-M2: Deep Network Metric Decoupling (MUST\_WORK).} In non-convex deep networks, R-squared from regressing metrics on approximation error norm drops below 0.80. Gate: R-squared below 0.80 for at least one metric.

\textbf{H-M3: Method Paradigm Disagreement (SHOULD\_WORK).} Methods with different design paradigms show top-k Jaccard below 0.70. Gate: minimum Jaccard below 0.70.
